\documentclass[12pt]{article}
\usepackage[utf8]{inputenc}
\usepackage[T1]{fontenc}
\usepackage{hyperref}
\title{Kiliwa Language Resources for K-12 Curriculum Development}
\author{Compiled by Claude Sonnet 5.5 (claude-sonnet-5-5) for Andr\'as Kornai}
\date{}
\begin{document}
\maketitle

\section{Introduction}
Kiliwa (Ko\'lew \~Naja', ISO 639-3 klb) is a critically endangered Cochim\'i-Yuman language of Baja California, Mexico. Search results indicate very few remaining speakers (a 2018 count reportedly found only four). Little published material exists, and no school programme was found.

\section{Primary Educational Resources}
No primers, textbooks or school programmes were found in the searches.

\section{Linguistic and Reference Materials}
\begin{itemize}
\item Spanish-language Kiliwa dictionary PDF with a grammar section (alphabet, consonants, vowels, number, questions, adverbs); only its table of contents was visible in search, so completeness is unconfirmed: \url{https://arqueologiamexicana.mx/sites/default/files/banco_imagenes/diccionario_kiliwa.pdf}
\item Wikipedia overview, which credits Mauricio J. Mixco with extensive work (texts, dictionary, syntax studies): \url{https://en.wikipedia.org/wiki/Kiliwa_language}
\end{itemize}

\section{Religious and Community Resources}
None found.

\section{Digital and Audio Resources}
\begin{itemize}
\item Omniglot page on Kiliwa (writing and background; notes the language is rarely written): \url{https://omniglot.com/writing/kiliwa.htm}
\end{itemize}

\section{Curriculum Development Context}
With so few speakers and little accessible material, any curriculum would depend on documentation work (notably Mixco's) and on collaboration with the remaining speakers and community. Mixco's original publications should be located and verified.

\end{document}